\section{Introduction}
\label{sec:intro}

Multimodal Large Language Models (MLLMs)~\cite{alayrac2022flamingo, liu2023visual, liu2023improved, bai2025qwen2} unify tasks involving multiple modalities, such as text, images and videos~\cite{caffagni2024r,sarto2025image}. Many of these tasks can be framed as Visual Question Answering (VQA)~\cite{goyal2017making,antol2015vqa,wu2017visual}, where a query may require understanding visual content, and the model must generate a faithful, correctly formatted response. Despite their broad pre-training corpora, even state-of-the-art MLLMs struggle with underrepresented, domain-specific queries~\cite{chen2023can, mensink2023encyclopedic}. This problem, known as Knowledge-based Visual Question Answering (KB-VQA)~\cite{marino2019ok}, is commonly addressed by enriching MLLMs with domain-specific information from external sources via \ie, Retrieval-Augmented Generation (RAG)~\cite{lewis2020retrieval}.

\begin{figure}[t]
    \centering
    \includegraphics[width=0.98\linewidth]{images/first-grpo.pdf}
    \vspace{-0.1cm}
    \caption{Comparison between Zero-Shot (ZS) MLLMs, retrieval-augmented models, and \ours. ZS MLLMs lack specialized knowledge and fail on domain-specific queries (top). Retrieval-augmented models introduce external context but often add noisy or irrelevant passages (middle). \ours overcomes this with a filtering stage over retrieved content and a multi-stage training strategy to enhance reasoning over passages.
    }
    \vspace{-0.3cm}
    \label{fig:first}
\end{figure}

Despite impressive results~\cite{yan2024echosight, wu2025towards, zhang2024mr, hong2025knowledge, cocchi2025augmenting, yuan2025mkg}, this setting still presents significant open challenges. One lies in information retrieval itself, as users’ queries can be extremely heterogeneous while the external knowledge-base can reach millions of documents in cardinality~\cite{chen2023can, mensink2023encyclopedic} -- thus lowering the recall of the retrieved results and adding noise to the MLLM input. This is further exacerbated by feature extraction and integration issues when queries and documents are multimodal, which is rapidly becoming usual~\cite{wei2023uniir, lin2024preflmr,caffagni2025recurrence}. Next, even assuming that the retrieved documents were relevant to the query, understanding them and extracting the right piece of information to generate the answer is not trivial~\cite{hua2025vision}.


To address these challenges, we propose \textbf{\ours}, short for \underline{\textbf{Re}}asoning-\underline{\textbf{A}}ugmented \underline{\textbf{G}}eneration, a novel multimodal retrieval-augmented generation approach that (i) mitigates low-recall and noisy retrieval by employing a multi-level retrieval pipeline followed by a critic model that effectively filters out irrelevant samples, and (ii) equips the MLLM with the capability of reasoning over retrieved results through a dedicated reinforcement learning training protocol. 

During retrieval, 
following common practice, \ours first employs an off-the-shelf multimodal encoder~\cite{radford2021learning, sun2024eva} for coarse-grained embedding-based retrieval, which achieves high recall when retaining many documents but suffers from low precision due to noise. To improve precision, a fine-grained retrieval stage focuses on the visual regions most relevant to the question. A critic model then classifies each passage as \textit{relevant} or \textit{irrelevant}, ensuring that only high-quality documents are passed to the generator (Fig.~\ref{fig:first}).


Building on advances in reasoning models~\cite{guo2025deepseek}, \ours further improves generation quality by allowing the model to produce explicit natural-language reasoning traces before the final answer.
Unlike prior multimodal RAG methods that rely on implicit reflection signals~\cite{asaiself, cocchi2025augmenting}, we train the model to create a full reasoning trace in natural language, without obeying any predefined pattern. To achieve this, \ours replaces the traditional supervised training in favor of a reinforcement learning framework inspired by GRPO~\cite{shao2024deepseekmath,yu2025dapo}, equipped with a reward scheme tailored for KB-VQA, to refine the model ability to reason over the user query and retrieved evidence, while using supervised fine-tuning only as a cold-start to establish initial reasoning behavior.


Experimentally, we evaluate the proposed approach on Encyclopedic-VQA~\cite{mensink2023encyclopedic} and InfoSeek~\cite{chen2023can}, two benchmarks containing question-answer pairs linked to Wikipedia-derived knowledge bases. Extensive experiments show that \ours substantially outperforms prior methods, not only improving answer accuracy but also generating explicit reasoning traces. These traces provide insight into the usefulness of retrieved passages and the steps leading to the final answer, offering full explainability of the model predictions.

\noindent In summary, the contributions of this work are as follows:
\begin{itemize}
    \item We propose \ours, a novel reasoning-augmented multimodal RAG model that combines coarse- and fine-grained retrieval with a critic model to improve precision and reduce noise injected to the generator. Notably, the critic is agnostic to the retrieval backbone, making it seamlessly applicable on top of any state-of-the-art retrieval engine.
   \item \ours trains the generator employing a multi-stage training strategy, leveraging SFT only as a cold start, followed by a reinforcement learning framework inspired by GRPO, with a reward scheme specifically designed for KB-VQA.
    \item We empirically validate \ours on two popular and challenging KB-VQA benchmarks, Encyclopedic-VQA and InfoSeek, where \ours reaches a new state-of-the-art.
\end{itemize}
